% Tabla D-T7. Generada automáticamente; no editar a mano.
% Requiere \usepackage{booktabs} (y graphicx si la tabla es ancha).
\begin{table}[htbp]
    \centering
    \caption{Verificación de los parámetros derivados frente a los valores del catálogo.}
    \label{tab:D-T7}
    \small
    \resizebox{\textwidth}{!}{%
    \begin{tabular}{lrrrrr}
    \toprule
    Comparación & $N$ & Razón mediana & $Q_1$ & $Q_3$ & $|\log_{10}$ razón$|$ mediana \\
    \midrule
    Semieje mayor: catalogado / Kepler & 5\,513 & 1.000 & 0.994 & 1.006 & 0.002 \\
    Flujo incidente: catalogado / calculado & 5\,644 & 0.997 & 0.919 & 1.080 & 0.035 \\
    \bottomrule
    \end{tabular}
    }
    \par\smallskip
    {\footnotesize\textit{Nota:} El flujo incidente se calculó para 5948 planetas; en 412 de ellos el semieje mayor se derivó con la tercera ley de Kepler.}
\end{table}
